\subsection{Calculus Vol2 - 7.2 Calculus of Parametric Curves}

\url{https://openstax.org/books/calculus-volume-2/pages/7-2-calculus-of-parametric-curves}

\begin{enumerate}
  \item How do you compute the derivative $\frac{dy}{dx}$ for a parametric curve defined by $x(t)$ and $y(t)$?
  \item Why is the formula $\frac{dy}{dx} = \frac{dy/dt}{dx/dt}$ valid, and what assumption is required about $\frac{dx}{dt}$?
  \item How do you find the equation of the tangent line to a parametric curve at a specific value $t=t_0$?
  \item How can you determine where a parametric curve has a horizontal tangent?
  \item How can you determine where a parametric curve has a vertical tangent?
  \item What happens to $\frac{dy}{dx}$ when $\frac{dx}{dt} = 0$?
  \item How do you compute the second derivative $\frac{d^2 y}{dx^2}$ for a parametric curve?
  \item Why is the second derivative useful when analyzing concavity of parametric curves?
  \item How can you interpret parametric equations as describing motion of a particle?
  \item What are the velocity components of a particle moving along a parametric curve?
  \item How is speed related to $x'(t)$ and $y'(t)$?
  \item How do you compute the arc length of a parametric curve between $t=a$ and $t=b$?
  \item What is the formula $L = \int_a^b \sqrt{(x'(t))^2 + (y'(t))^2},dt$ and how is it derived conceptually?
  \item How do you compute the area under a parametric curve?
  \item Why is the formula $A = \int_{t_1}^{t_2} y(t),x'(t),dt$ used for area?
  \item How does the sign of $x'(t)$ affect the computed area?
  \item What conditions must hold for the area formula for parametric curves to be valid?
  \item How do you handle curves that are traced more than once when computing area?
  \item How can you determine whether a parametric curve is traced once or multiple times?
  \item How do you compute the surface area generated by revolving a parametric curve about the $x$-axis?
  \item What is the formula $S = 2\pi \int_a^b y(t)\sqrt{(x'(t))^2 + (y'(t))^2},dt$ and when is it used?
  \item How does the surface area formula change when revolving around the $y$-axis?
  \item What geometric idea underlies the surface area of revolution for parametric curves?
\end{enumerate}
